% Ordered Associative Chiral Koszul Duality
% Integration-ready chapter file (stripped from standalone amsart draft).
% Uses only \providecommand for macros that may not be in main.tex preamble.

\providecommand{\Assch}{\mathrm{Ass}^{\mathrm{ch}}}
\providecommand{\Barch}{\overline{B}^{\mathrm{ch}}}
\providecommand{\Cobar}{\Omega^{\mathrm{ch}}}
\providecommand{\coHoch}{\operatorname{coHH}}
\providecommand{\Cotor}{\operatorname{Cotor}}
\providecommand{\Coext}{\operatorname{Coext}}
\providecommand{\RHom}{R\!\operatorname{Hom}}
\providecommand{\Tot}{\operatorname{Tot}}
\providecommand{\KK}{\mathbb{K}}
\providecommand{\Dpbw}{D^{\mathrm{pbw}}}
\providecommand{\Dco}{D^{\mathrm{co}}}
\providecommand{\chotimes}{\mathbin{\otimes^{\mathrm{ch}}}}
\providecommand{\wt}{\widetilde}
\providecommand{\eps}{\varepsilon}
\providecommand{\susp}{s}
\providecommand{\coeq}{\operatorname{coeq}}
\providecommand{\id}{\mathrm{id}}
\providecommand{\op}{\mathrm{op}}
\providecommand{\cop}{\mathrm{cop}}
\providecommand{\HH}{\operatorname{HH}}
\providecommand{\Tor}{\operatorname{Tor}}
\providecommand{\Ext}{\operatorname{Ext}}
\providecommand{\gr}{\operatorname{gr}}
\providecommand{\MC}{\operatorname{MC}}
\providecommand{\Ob}{\operatorname{Ob}}
\providecommand{\eq}{\operatorname{eq}}

\chapter{Ordered Associative Chiral Koszul Duality}
\label{ch:ordered-associative-chiral-kd}

\section{Introduction: the diagonal as the universal boundary}

The commutative shadow of chiral Koszul duality is by now classical.  The genuinely associative
story is different.  It is not the Francis--Gaitsgory equivalence between chiral Lie algebras and
cocommutative factorization coalgebras.  It is a theory in which ordered configuration spaces
survive, shuffles do not collapse under symmetric-group quotients, and after linear duality
chiral algebras appear again on the far side.

The purpose of this note is to make the associative picture rigid enough to be useful.  We do this
by returning to first principles and asking a single question.

\medskip

\emph{What is the correct coalgebra-side avatar of the diagonal bimodule ${}_A A_A$?}

\medskip

The answer is simple and structural:
if $C=\Barch(A)$ is the ordered bar coalgebra of an associative chiral algebra $A$, then the
correct boundary object is not an ad hoc kernel but the \emph{diagonal bicomodule}
\[
C_\Delta,
\]
namely the coalgebra $C$ equipped with its coproduct on the left and on the right.  Once this is
recognized, the rest of the machine begins to lock into place.

On the algebra side, interval composition is derived tensor over $A$ and the annulus is Hochschild
homology.  On the coalgebra side, interval composition becomes derived cotensor over $C$, the
diagonal becomes the composition unit, and the annulus becomes the derived cotrace of the diagonal
against itself:
\[
\HH^{\mathrm{ch}}_\bullet(A)\simeq C_\Delta \square_{C^e}^{\mathbf R} C_\Delta.
\]
In other words, the annulus is no longer mysterious.  It is the self-cotensor of the universal
boundary condition.

This note is organized as follows.  In \S\ref{sec:setup} we isolate the exact regime in which the
homological algebra is clean and state the repository inputs used below.  In \S\ref{sec:shuffle} we
prove the ordered chiral shuffle theorem.  In \S\ref{sec:opposite} we prove opposite-duality and
enveloping duality.  In \S\ref{sec:tangent} we construct the two-sided twisted bicomodule and prove
that it is exactly the one-defect ordered bar complex of a square-zero extension.  In
\S\ref{sec:bimod-bicomod} we obtain the PBW-complete equivalence between bimodules and bicomodules
and prove that ${}_A A_A$ goes to $C_\Delta$.  In \S\ref{sec:HH-coHH} we identify associative chiral
Hochschild homology and cohomology with coalgebraic coHochschild theory.  In
\S\ref{sec:open-trace} we extract the ordered genus-zero open trace formalism and its Calabi--Yau
refinement.  Finally, in \S\ref{sec:conjectures} we formulate the next frontier as a precise list of
conjectures.

\subsection*{Dictionary}
The new results of this note can be summarized by the following dictionary:
\[
A \longleftrightarrow C_\Delta,\qquad
A^e \longleftrightarrow C^e:=C\widehat\otimes C^{\cop},\qquad
-\otimes_A^{\mathbf L}- \longleftrightarrow -\square_C^{\mathbf R}-,
\]
\[
\HH^{\mathrm{ch}}_\bullet(A,M)
\longleftrightarrow
\coHoch^{\mathrm{ch}}_\bullet(C,\KK_A^{\mathrm{bi}}(M)),
\qquad
\HH_{\mathrm{ch}}^\bullet(A,M)
\longleftrightarrow
\coHoch_{\mathrm{ch}}^\bullet(C,\KK_A^{\mathrm{bi}}(M)).
\]
The slogan is: \emph{the associative dual does not merely mirror the algebra; it already remembers
the open boundary formalism.}

\subsection*{Ledger of status}
To keep the exposition honest, we separate three classes of statements.

\begin{enumerate}[label=(\roman*)]
\item \textbf{Repository inputs.}
The existence of the associative chiral bar--cobar equivalence, the one-sided module/comodule
Koszul duality, and the enveloping algebra formalism are taken as already established in the
repository.

\item \textbf{New theorems proved here.}
The ordered chiral shuffle theorem, opposite-duality, enveloping duality, the two-sided twisted
bicomodule, the identification of the tangent bar complex with the twisted bicomodule, the PBW-complete
bimodule/bicomodule equivalence, the diagonal correspondence, the Hochschild--coHochschild
dictionary, and the ordered open trace formalism.

\item \textbf{Conjectural frontier.}
The bordered configuration-space realization, the Francis--Gaitsgory commutator-shadow spectral
sequence, the associative modular Maurer--Cartan class, and BRST/Drinfeld--Sokolov compatibility.
\end{enumerate}

\section{The strongly admissible regime and the foundational inputs}
\label{sec:setup}

We work over a field $k$ of characteristic $0$.  Let $X$ be a smooth curve over $k$, and let
$\mathcal D(X)$ denote the relevant derived category of $D$-modules on $X$ (or on the Ran space of
$X$) endowed with the chiral tensor product
\[
\chotimes := \boxtimes_{D_X}.
\]
An $E_1$-chiral algebra means an algebra over the non-$\Sigma$ operad $\Assch$.

\begin{definition}[Strongly admissible associative chiral algebra]
An augmented $E_1$-chiral algebra $A$ on $X$ is called \emph{strongly admissible} if the following
hold.
\begin{enumerate}[label=(\alph*)]
\item $A$ is pro-nilpotent and equipped with an exhaustive separated descending PBW-type filtration
$F^\bullet A$ which is complete for the induced topology.

\item Every associated graded piece and every finite filtration quotient is bounded below and perfect
enough that completed tensor products, totalizations, and normalized chains behave exact\-ly on these
pieces.

\item The reduced ordered bar coalgebra
\[
C:=\Barch(A)
\]
is conilpotent, and the corresponding coderived category of bicomodules admits enough coflat
resolutions so that derived cotensor products are defined and associative.

\item Whenever linear duality is invoked, each graded/conformal piece of the object under consideration
is finite-dimensional.
\end{enumerate}
\end{definition}

\begin{remark}
This is the fortified regime in which hidden homological-algebra issues are kept on the surface.
Nothing deep in the new arguments below requires more than these hypotheses; they are chosen so that
every totalization, completion, normalization, and passage to coderived categories is legitimate.
\end{remark}

\begin{hypothesis}[Repository inputs]
\label{hyp:inputs}
We assume the repository has already established the following.
\begin{enumerate}[label=(I)]
\item \textbf{Associative chiral bar--cobar.}
For every augmented pro-nilpotent strongly admissible $E_1$-chiral algebra $A$ there is a
quasi-isomorphism
\[
\Cobar \Barch(A)\xrightarrow{\sim} A,
\]
and, on the appropriate coalgebra side, the corresponding bar--cobar adjunction is an equivalence.

\item \textbf{One-sided module/comodule duality.}
For every such $A$ there is an equivalence between complete left $A$-modules and conilpotent left
$\Barch(A)$-comodules, with the usual coderived/contraderived refinement in the curved or completed
regime.

\item \textbf{Enveloping algebra formalism.}
The chiral enveloping algebra
\[
A^e := A\chotimes A^{\op}
\]
is defined, and $A$-bimodules are equivalently left $A^e$-modules.
\end{enumerate}
\end{hypothesis}

\begin{definition}[Universal twisting cochain]
Let $A$ be strongly admissible and set $C:=\Barch(A)$.  The bar--cobar equivalence produces a
canonical twisting cochain
\[
\tau:C\to A
\]
characterized by the Maurer--Cartan equation
\[
d\tau+\tau\star\tau=0.
\]
We call $\tau$ the \emph{universal twisting cochain}.
\end{definition}

\begin{definition}[Opposite and diagonal objects]
Let $A$ be an associative chiral algebra and let $C$ be a coaugmented dg coalgebra.
\begin{enumerate}[label=(\alph*)]
\item $A^{\op}$ is the same underlying object with multiplication reversed.
\item $C^{\cop}$ is the same underlying object with coproduct flipped.
\item The \emph{coalgebraic enveloping object} is
\[
C^e:=C\widehat\otimes C^{\cop}.
\]
\item The \emph{diagonal bicomodule} $C_\Delta$ is $C$ equipped with left and right coactions both
equal to $\Delta_C$.
\end{enumerate}
\end{definition}

\begin{lemma}[Bicomodules as comodules over the enveloping coalgebra; \ClaimStatusProvedHere]
\label{lem:bicom-e}
Let $C$ be a coaugmented conilpotent dg coalgebra.  Then left $C^e$-comodules are canonically the
same thing as $C$-bicomodules.
\end{lemma}

\begin{proof}
Given a left $C^e$-coaction
\[
\rho^e:M\to C\otimes C^{\cop}\otimes M,
\]
compose with the counit on the second factor and on the first factor to obtain
\[
\lambda=(\id_C\otimes \eps\otimes \id_M)\rho^e:M\to C\otimes M,
\]
\[
\rho=(\eps\otimes \id_{C^{\cop}}\otimes \id_M)\rho^e:M\to M\otimes C.
\]
These are a left and a right $C$-coaction, and coassociativity of $\rho^e$ is equivalent to the
usual bicomodule compatibility condition.

Conversely, given commuting left and right coactions $\lambda$ and $\rho$, define
\[
\rho^e
=
(\id_C\otimes \tau_{C,M}\otimes \id_M)\circ(\lambda\otimes \id_M)\circ \rho
:
M\to C\otimes C^{\cop}\otimes M,
\]
where $\tau_{C,M}$ swaps the middle two factors.  The bicomodule axioms are precisely the
coassociativity and counitality conditions for $\rho^e$.
\end{proof}

\begin{definition}[Chiral Hochschild and coHochschild theories]
Let $A$ be a strongly admissible associative chiral algebra and $C=\Barch(A)$.
For an $A$-bimodule $M$ and a $C$-bicomodule $N$, define
\[
\HH^{\mathrm{ch}}_\bullet(A,M):=\Tor_\bullet^{A^e}(A,M),
\qquad
\HH_{\mathrm{ch}}^\bullet(A,M):=\RHom_{A^e}(A,M),
\]
\[
\coHoch^{\mathrm{ch}}_\bullet(C,N):=\Cotor_\bullet^{C^e}(C_\Delta,N),
\qquad
\coHoch_{\mathrm{ch}}^\bullet(C,N):=\Coext_{C^e}^\bullet(C_\Delta,N).
\]
When $M=A$ and $N=C_\Delta$ we omit coefficients.
\end{definition}

\section{Ordered bar coalgebras and the chiral shuffle theorem}
\label{sec:shuffle}

In the associative chiral regime the basic bar object is ordered.  This is the first point at which
the theory diverges sharply from the commutative shadow: \emph{every shuffle survives as a distinct
summand.}

\begin{definition}[The simplicial ordered bar object]
Let $A$ be an augmented associative chiral algebra.  Define the simplicial object
$B^\Delta_\bullet(A)$ by
\[
B^\Delta_n(A):=A^{\chotimes n}\qquad (n\geq 0),
\]
with faces given by augmentation at the two ends and multiplication in adjacent slots, and
degeneracies given by unit insertion.  Its normalized chains
\[
N(B^\Delta_\bullet(A))
\]
are canonically identified with the reduced ordered chiral bar complex $\Barch(A)$.
\end{definition}

In suspended notation we write a bar word of length $n$ as
\[
[a_1|\cdots|a_n].
\]

\begin{definition}[EZ-admissible pair]
A pair $(A,B)$ of strongly admissible associative chiral algebras is called \emph{EZ-admissible} if
the normalization functor on bisimplicial objects built from $A$ and $B$ commutes, filtration-wise,
with the completed chiral tensor products and totalizations appearing in the proof of the
Eilenberg--Zilber theorem.
\end{definition}

\begin{theorem}[Ordered chiral shuffle theorem; \ClaimStatusProvedHere]
\label{thm:shuffle}
Let $(A,B)$ be an EZ-admissible pair of strongly admissible associative chiral algebras.
There are natural filtered chain maps
\[
\mathsf{Sh}^{\mathrm{ch}}_{A,B}:
\Barch(A)\widehat\otimes \Barch(B)\longrightarrow \Barch(A\chotimes B),
\]
\[
\mathsf{AW}^{\mathrm{ch}}_{A,B}:
\Barch(A\chotimes B)\longrightarrow \Barch(A)\widehat\otimes \Barch(B),
\]
such that
\[
\mathsf{AW}^{\mathrm{ch}}_{A,B}\circ \mathsf{Sh}^{\mathrm{ch}}_{A,B}=\id,
\qquad
\mathsf{Sh}^{\mathrm{ch}}_{A,B}\circ \mathsf{AW}^{\mathrm{ch}}_{A,B}\simeq \id.
\]
Hence there is a filtered quasi-isomorphism of conilpotent dg coalgebras
\[
\Barch(A)\widehat\otimes \Barch(B)\xrightarrow{\sim}\Barch(A\chotimes B).
\]
\end{theorem}

\begin{proof}
Set
\[
K_{\bullet,\bullet}:=B^\Delta_\bullet(A)\chotimes B^\Delta_\bullet(B).
\]
This is a bisimplicial object in the ambient abelian or stable category of $D$-modules.

Its diagonal identifies canonically with the simplicial bar object of the tensor-product algebra:
\[
\operatorname{diag}(K)_n
=
A^{\chotimes n}\chotimes B^{\chotimes n}
\cong
(A\chotimes B)^{\chotimes n}
=
B^\Delta_n(A\chotimes B).
\]
This identification is compatible with faces and degeneracies because multiplication and augmentation
on $A\chotimes B$ are componentwise.

Now apply the classical Eilenberg--Zilber theorem filtration-wise in each exact graded piece.
EZ-admissibility ensures that normalization commutes with the completed tensor constructions used
here, so one obtains a filtered quasi-isomorphism
\[
N(\operatorname{diag}K)\xrightarrow{\sim} \Tot\bigl(NK_{\bullet,\bullet}\bigr),
\]
with inverse given by Alexander--Whitney.  Since
\[
N(B^\Delta_\bullet(A))\cong \Barch(A),
\qquad
N(B^\Delta_\bullet(B))\cong \Barch(B),
\qquad
N(B^\Delta_\bullet(A\chotimes B))\cong \Barch(A\chotimes B),
\]
this yields the claimed quasi-isomorphism.

The formulas are the usual normalized Eilenberg--MacLane formulas, now interpreted in the ordered
chiral setting.  Explicitly, for bar words
\[
[a_1|\cdots|a_p]\in \Barch_p(A),\qquad [b_1|\cdots|b_q]\in \Barch_q(B),
\]
set $c_i=a_i\otimes 1$ for $1\leq i\leq p$ and $c_{p+j}=1\otimes b_j$ for $1\leq j\leq q$, and define
\[
\mathsf{Sh}^{\mathrm{ch}}_{p,q}
([a_1|\cdots|a_p]\otimes[b_1|\cdots|b_q])
=
\sum_{\sigma\in \operatorname{Sh}(p,q)}
\varepsilon(\sigma)\,
[c_{\sigma^{-1}(1)}|\cdots|c_{\sigma^{-1}(p+q)}].
\]
Similarly,
\[
\mathsf{AW}^{\mathrm{ch}}([x_1|\cdots|x_n])
=
\sum_{p=0}^n
[\pi_A(x_1)|\cdots|\pi_A(x_p)]\otimes[\pi_B(x_{p+1})|\cdots|\pi_B(x_n)],
\]
where
\[
\pi_A=\id_A\otimes \eps_B,\qquad \pi_B=\eps_A\otimes \id_B.
\]
The chain-map property uses only the simplicial identities.

Finally, the shuffle map is compatible with deconcatenation: cutting a shuffled word is equivalent to
cutting the source words and then shuffling the left halves and the right halves.  Therefore
$\mathsf{Sh}^{\mathrm{ch}}_{A,B}$ and $\mathsf{AW}^{\mathrm{ch}}_{A,B}$ are coalgebra maps.
\end{proof}

\begin{remark}
The theorem is genuinely associative.  In a commutative or $E_\infty$ setting many of the shuffle
summands are identified after the symmetric-group quotient.  Here they remain distinct.  The ordered
configuration space remembers their full combinatorics.
\end{remark}

\section{Opposite-duality, anti-involutions, and enveloping duality}
\label{sec:opposite}

The second foundational phenomenon is order reversal.  It transforms opposite multiplication on the
algebra side into opposite deconcatenation on the bar side.

\begin{theorem}[Opposite-duality for ordered bar coalgebras; \ClaimStatusProvedHere]
\label{thm:opposite}
Let $A$ be a strongly admissible associative chiral algebra.  There is a canonical isomorphism of
coaugmented conilpotent dg coalgebras
\[
\mathsf{R}_A:\Barch(A^{\op})\xrightarrow{\sim}\Barch(A)^{\cop}.
\]
If $A$ is finite-type Koszul so that linear duality returns an associative chiral algebra, then
\[
(A^{\op})^!\simeq (A^!)^{\op}.
\]
\end{theorem}

\begin{proof}
On suspended bar words define
\[
\mathsf{R}_A[\susp a_1|\cdots|\susp a_n]
=
(-1)^{\sum_{i<j}(|a_i|+1)(|a_j|+1)}
[\susp a_n|\cdots|\susp a_1].
\]
Geometrically, this is induced by the order-reversal involution on the ordered configuration space.

The bar differential is the alternating sum of adjacent collision maps.  The $i$-th face for
$A^{\op}$ multiplies the $i$-th and $(i+1)$-st slots in the opposite order.  Reversal carries that
face to the $(n-i)$-th face for $A$ with exactly the sign written above.  Hence $\mathsf{R}_A$ is a
chain map.

The coproduct on $\Barch(A)$ is deconcatenation.  Reversal sends a cut after the $p$-th letter to a
cut before the $(n-p)$-th letter.  Therefore
\[
\Delta\circ\mathsf{R}_A
=
\tau\circ(\mathsf{R}_A\otimes\mathsf{R}_A)\circ\Delta,
\]
where $\tau$ flips the tensor factors, which is exactly the statement that the target is the
opposite coalgebra.  The claim about linear duals follows because linear duality exchanges opposite
coproduct with opposite product.
\end{proof}

\begin{corollary}[Anti-involutions survive duality; \ClaimStatusProvedHere]
\label{cor:anti}
Any anti-involution $\iota:A\to A^{\op}$ induces an anti-involution
\[
\iota^!:A^!\to (A^!)^{\op}.
\]
If $\iota^2=\id$, then $(\iota^!)^2=\id$.
\end{corollary}

\begin{proof}
Apply Theorem~\ref{thm:opposite} functorially.
\end{proof}

\begin{lemma}[Closure of admissibility under opposite and enveloping constructions; \ClaimStatusProvedHere]
\label{lem:closure}
If $A$ is strongly admissible, then so are $A^{\op}$ and $A^e=A\chotimes A^{\op}$.
\end{lemma}

\begin{proof}
All hypotheses in the definition of strong admissibility are stable under replacing multiplication by
its opposite and under taking completed tensor products of filtered objects.
\end{proof}

\begin{corollary}[Enveloping duality; \ClaimStatusProvedHere]
\label{cor:enveloping}
Let $A$ be strongly admissible, and assume $(A,A^{\op})$ is EZ-admissible.  Set $C=\Barch(A)$.
Then there is a canonical quasi-isomorphism of conilpotent dg coalgebras
\[
\Barch(A^e)\xrightarrow{\sim} C^e=C\widehat\otimes C^{\cop}.
\]
If $A$ is finite-type Koszul, then
\[
(A^e)^!\simeq (A^!)^e.
\]
\end{corollary}

\begin{proof}
By the ordered shuffle theorem,
\[
\Barch(A^e)=\Barch(A\chotimes A^{\op})\xrightarrow{\sim}\Barch(A)\widehat\otimes \Barch(A^{\op}).
\]
Now apply Theorem~\ref{thm:opposite} to the second factor:
\[
\Barch(A^{\op})\xrightarrow{\sim}\Barch(A)^{\cop}.
\]
This proves the coalgebra statement.  The algebra statement follows after finite-type linear duality.
\end{proof}

\section{The two-sided twisted bicomodule and the tangent bar construction}
\label{sec:tangent}

We now build the correct coalgebra-side avatar of an $A$-bimodule.  There are two a priori
different constructions:
the first is obtained by differentiating the ordered bar construction at a square-zero extension;
the second is obtained by twisting a two-sided tensor product with the universal twisting cochain.
A central point of this note is that these two constructions are the same.

\begin{definition}[Two-sided twisted bicomodule]
\label{def:Kbi}
Let $A$ be strongly admissible, let $C=\Barch(A)$, and let $M$ be an $A$-bimodule.
Define the \emph{two-sided twisted bicomodule}
\[
\KK_A^{\mathrm{bi}}(M):=C\widehat\otimes M\widehat\otimes C
\]
with differential
\[
d_{\KK}
=
d_C\otimes 1\otimes 1 + 1\otimes d_M\otimes 1 + 1\otimes 1\otimes d_C + d_L+d_R,
\]
where
\[
d_L(c\otimes m\otimes d)
=
\sum (-1)^{|c_{(1)}|}\, c_{(1)}\otimes \tau(c_{(2)})\cdot m\otimes d,
\]
\[
d_R(c\otimes m\otimes d)
=
\sum (-1)^{|c|+|m|}\, c\otimes m\cdot \tau(d_{(1)})\otimes d_{(2)},
\]
with the usual Koszul signs determined by the grading conventions.
The left and right $C$-coactions are induced by comultiplication on the two outer factors.
\end{definition}

\begin{lemma}[\ClaimStatusProvedHere]
\label{lem:Kbi-dg}
$\KK_A^{\mathrm{bi}}(M)$ is a dg $C$-bicomodule.
\end{lemma}

\begin{proof}
The untwisted part of the differential squares to zero.  The quadratic terms in $d_L$ and in $d_R$
vanish by the Maurer--Cartan equation
\[
d\tau+\tau\star\tau=0.
\]
The mixed commutator $d_Ld_R+d_Rd_L$ vanishes because the left and right $A$-actions on the
bimodule $M$ commute.  Compatibility with the bicomodule structure is immediate from coassociativity
of the coproduct on the outer $C$-factors.
\end{proof}

\begin{definition}[Square-zero extension and one-defect bar complex]
Let $M$ be an $A$-bimodule and form the square-zero extension
\[
A_\eps:=A\oplus \eps M,\qquad \eps^2=0.
\]
Inside the reduced ordered bar complex $\Barch(A_\eps)$, let $T_A^{\mathrm{bar}}(M)$ be the
subcomplex spanned by bar words containing exactly one $\eps M$-slot.
We call this the \emph{one-defect ordered bar complex}.
\end{definition}

\begin{proposition}[\ClaimStatusProvedHere]
\label{prop:one-defect}
Modulo $\eps^2$, there is a canonical decomposition
\[
\Barch(A_\eps)\cong \Barch(A)\oplus \eps\, T_A^{\mathrm{bar}}(M).
\]
Moreover, $T_A^{\mathrm{bar}}(M)$ is naturally a $C$-bicomodule.
\end{proposition}

\begin{proof}
A reduced bar word in $\Barch(A_\eps)$ contains some number of $\eps M$-slots.
Since $\eps^2=0$, any word with at least two such slots vanishes modulo $\eps^2$.
Therefore only the zero-defect and one-defect pieces survive.

The bar differential merges adjacent slots.  Multiplying two ordinary $A$-slots remains in $A$.
Multiplying an $A$-slot with the unique $M$-slot remains in $M$ by the bimodule structure.
There is no nonzero product of two defect slots.  Hence the one-defect sector is a subcomplex.
Deconcatenation cuts a one-defect word into a defect-free part and a one-defect part, or vice versa,
so the one-defect sector is naturally a bicomodule over the ordinary bar coalgebra $C=\Barch(A)$.
\end{proof}

\begin{theorem}[Tangent identification; \ClaimStatusProvedHere]
\label{thm:tangent=K}
There is a canonical isomorphism of dg $C$-bicomodules
\[
T_A^{\mathrm{bar}}(M)\xrightarrow{\sim}\KK_A^{\mathrm{bi}}(M).
\]
Equivalently, the one-defect ordered bar complex of a square-zero extension is exactly the two-sided
twisted bicomodule of the defect bimodule.
\end{theorem}

\begin{proof}
As a graded object,
\[
C=\bigoplus_{p\geq 0} (\susp\bar A)^{\otimes p},
\]
hence
\[
C\widehat\otimes M\widehat\otimes C
\cong
\bigoplus_{p,q\geq 0}
(\susp\bar A)^{\otimes p}\otimes M\otimes(\susp\bar A)^{\otimes q}.
\]
This is in tautological bijection with one-defect bar words
\[
[a_1|\cdots|a_p|m|a_{p+1}|\cdots|a_{p+q}]
\]
in which the unique defect sits between a left bar segment of length $p$ and a right bar segment of
length $q$.

Under this identification, the ordinary bar differential on the left segment and the right segment
matches the internal bar differential coming from the two outer copies of $C$, while the two bar
faces adjacent to the defect slot match exactly the two twisting terms $d_L$ and $d_R$.
Therefore the differentials agree.  Deconcatenation also agrees on the nose, because cutting a
one-defect word is the same as cutting the left and right bar segments around the unique defect.
\end{proof}

\begin{corollary}[Infinitesimal dual coalgebra; \ClaimStatusProvedHere]
\label{cor:infdual}
Modulo $\eps^2$, the bar coalgebra of the square-zero extension is
\[
\Barch(A_\eps)\cong C\oplus \eps\,\KK_A^{\mathrm{bi}}(M).
\]
If $A$ is finite-type Koszul, then after linear duality
\[
(A_\eps)^!\cong A^!\oplus \eps\, M^!
\qquad (\mathrm{mod}\ \eps^2),
\]
where
\[
M^!:=\Omega_C^{\mathrm{bi}}\bigl(\KK_A^{\mathrm{bi}}(M)\bigr).
\]
\end{corollary}

\begin{proof}
Combine Proposition~\ref{prop:one-defect} and Theorem~\ref{thm:tangent=K}.  The statement after
dualization is the corresponding first-order cobar reconstruction.
\end{proof}

\begin{proposition}[Infinitesimal annular variation; \ClaimStatusProvedHere]
\label{prop:infann}
Modulo $\eps^2$ one has
\[
\HH^{\mathrm{ch}}_\bullet(A_\eps)
\cong
\HH^{\mathrm{ch}}_\bullet(A)\oplus \eps\, \HH^{\mathrm{ch}}_\bullet(A,M).
\]
On the coalgebra side this becomes
\[
\coHoch^{\mathrm{ch}}_\bullet(\Barch(A_\eps))
\cong
\coHoch^{\mathrm{ch}}_\bullet(C)\oplus \eps\,
\coHoch^{\mathrm{ch}}_\bullet(C,\KK_A^{\mathrm{bi}}(M)).
\]
\end{proposition}

\begin{proof}
The Hochschild complex of the square-zero extension decomposes, modulo $\eps^2$, into the sector with
no defect and the sector with exactly one defect; the second computes Hochschild homology with
coefficients in $M$.  The coalgebra statement follows from
Corollary~\ref{cor:infdual} and the coHochschild identification proved later in
Theorem~\ref{thm:HH-coHH-homology}.
\end{proof}

\section{PBW-complete two-sided Koszul duality}
\label{sec:bimod-bicomod}

The repository already establishes one-sided module/comodule duality.  We now bootstrap it to the
two-sided theory by passing from $A$ to its enveloping algebra $A^e$ and then using
Corollary~\ref{cor:enveloping}.

\begin{definition}[PBW-complete and coderived bicomodule categories]
Let $A$ be strongly admissible and set $C=\Barch(A)$.
\begin{enumerate}[label=(\alph*)]
\item $\Dpbw_{\mathrm{bi}}(A)$ denotes the derived category of PBW-complete $A$-bimodules obtained
from complete semifree $A^e$-resolutions.

\item $\Dco_{\mathrm{bi}}(C)$ denotes the coderived category of conilpotent $C$-bicomodules.
Equivalently, by Lemma~\ref{lem:bicom-e}, it is the coderived category of left $C^e$-comodules.
\end{enumerate}
\end{definition}

\begin{definition}[Two-sided twisted cobar]
Let $N$ be a $C$-bicomodule.  Define
\[
\Omega_C^{\mathrm{bi}}(N):=A\widehat\otimes N\widehat\otimes A
\]
with the differential dual to the one in Definition~\ref{def:Kbi}.
\end{definition}

\begin{theorem}[PBW-complete bimodule/bicomodule equivalence; \ClaimStatusProvedHere]
\label{thm:bimod-bicomod}
Let $A$ be strongly admissible, and assume $(A,A^{\op})$ is EZ-admissible.
Set $C=\Barch(A)$.
Then there is an exact equivalence
\[
\KK_A^{\mathrm{bi}}:\Dpbw_{\mathrm{bi}}(A)\xrightarrow{\sim}\Dco_{\mathrm{bi}}(C)
\]
with inverse $\Omega_C^{\mathrm{bi}}$.  In the curved regime the same statement holds after replacing
the ordinary derived category on the algebra side by the corresponding completed
contra\-derived/coderived version.
\end{theorem}

\begin{proof}
By the enveloping algebra formalism, $A$-bimodules are left $A^e$-modules.
By Corollary~\ref{cor:enveloping},
\[
\Barch(A^e)\xrightarrow{\sim} C^e.
\]
Apply the one-sided module/comodule Koszul duality from Hypothesis~\ref{hyp:inputs}(II) to the
associative chiral algebra $A^e$.  This gives an equivalence between complete $A^e$-modules and
conilpotent $C^e$-comodules.  Since left $C^e$-comodules are the same as $C$-bicomodules by
Lemma~\ref{lem:bicom-e}, we obtain the asserted equivalence.

The functor is precisely the two-sided twisted bar construction
\[
M\longmapsto C\widehat\otimes M\widehat\otimes C,
\]
because that is the one-sided bar functor for the algebra $A^e$ rewritten using the enveloping
identification.  The inverse is the corresponding two-sided cobar functor.
\end{proof}

\begin{theorem}[Diagonal correspondence; \ClaimStatusProvedHere]
\label{thm:diagonal}
Under the equivalence of Theorem~\ref{thm:bimod-bicomod}, the diagonal bimodule ${}_A A_A$ is sent
to the diagonal bicomodule $C_\Delta$.  Equivalently,
\[
\KK_A^{\mathrm{bi}}(A)\simeq C_\Delta
\]
in $\Dco_{\mathrm{bi}}(C)$.
\end{theorem}

\begin{proof}
By definition,
\[
\KK_A^{\mathrm{bi}}(A)=C\widehat\otimes A\widehat\otimes C
\]
with the two-sided twisting differential.

First contract the left half:
the standard bar--cobar counit resolution yields a quasi-isomorphism
\[
C\widehat\otimes_\tau A\xrightarrow{\sim} \mathbf 1
\]
compatible with the right $C$-coaction.  Hence
\[
\KK_A^{\mathrm{bi}}(A)\simeq C
\]
as a right $C$-comodule.

Similarly, contracting the right half using
\[
A\widehat\otimes_\tau C\xrightarrow{\sim}\mathbf 1
\]
gives
\[
\KK_A^{\mathrm{bi}}(A)\simeq C
\]
as a left $C$-comodule.

To identify the opposite coactions, filter $\KK_A^{\mathrm{bi}}(A)$ by total coradical degree in the
two outer copies of $C$.  On the associated graded the twisting terms disappear, and the bicomodule
structure is literally the left and right coproducts on the two outer factors.  Under either of the
two contractions above, the associated graded collapses to a single copy of $C$, and the surviving
left and right coactions are exactly the two legs of $\Delta_C$.  Because the filtration is complete
and bounded below in the strongly admissible regime, this associated-graded identification lifts to
the filtered complex.  Therefore the resulting bicomodule is $C_\Delta$.
\end{proof}

\begin{corollary}[The diagonal is the unit for composition; \ClaimStatusProvedHere]
\label{cor:unit}
For every $X\in \Dco_{\mathrm{bi}}(C)$ there are canonical isomorphisms
\[
C_\Delta \square_C^{\mathbf R} X \simeq X \simeq X \square_C^{\mathbf R} C_\Delta.
\]
\end{corollary}

\begin{proof}
Transport the tautologies
\[
A\otimes_A^{\mathbf L} M\simeq M\simeq M\otimes_A^{\mathbf L} A
\]
for $A$-bimodules along Theorem~\ref{thm:bimod-bicomod} and use
Theorem~\ref{thm:diagonal}.
\end{proof}

\begin{corollary}[Tensor--cotensor gluing; \ClaimStatusProvedHere]
\label{cor:tensor-cotensor}
For all $M,N\in \Dpbw_{\mathrm{bi}}(A)$ there is a canonical isomorphism
\[
\KK_A^{\mathrm{bi}}(M\otimes_A^{\mathbf L} N)
\simeq
\KK_A^{\mathrm{bi}}(M)\square_C^{\mathbf R}\KK_A^{\mathrm{bi}}(N)
\]
in $\Dco_{\mathrm{bi}}(C)$.
\end{corollary}

\begin{proof}
Transport the derived tensor product across the equivalence of
Theorem~\ref{thm:bimod-bicomod}.
\end{proof}

\begin{remark}
Corollary~\ref{cor:tensor-cotensor} is the coalgebraic gluing law for intervals.
The diagonal bicomodule $C_\Delta$ is therefore the universal boundary condition for interval
composition.
\end{remark}

\section{Hochschild--coHochschild duality}
\label{sec:HH-coHH}

We now identify annular traces and centers on the algebra side with coHochschild theories of the bar
coalgebra.

\begin{theorem}[Associative chiral Hochschild/coHochschild homology; \ClaimStatusProvedHere]
\label{thm:HH-coHH-homology}
Let $A$ be strongly admissible, let $(A,A^{\op})$ be EZ-admissible, and set $C=\Barch(A)$.
Then for every complete $A$-bimodule $M$ there is a canonical isomorphism
\[
\HH^{\mathrm{ch}}_\bullet(A,M)\simeq
\coHoch^{\mathrm{ch}}_\bullet(C,\KK_A^{\mathrm{bi}}(M)).
\]
In particular,
\[
\HH^{\mathrm{ch}}_\bullet(A)\simeq \coHoch^{\mathrm{ch}}_\bullet(C)
\simeq
C_\Delta \square_{C^e}^{\mathbf R} C_\Delta.
\]
\end{theorem}

\begin{proof}
By definition,
\[
\HH^{\mathrm{ch}}_\bullet(A,M)=\Tor_\bullet^{A^e}(A,M).
\]
Apply the equivalence of Theorem~\ref{thm:bimod-bicomod} to the pair $(A,M)$.
The diagonal bimodule $A$ goes to $C_\Delta$ by Theorem~\ref{thm:diagonal}, and $M$ goes to
$\KK_A^{\mathrm{bi}}(M)$ by definition.  Derived tensor over $A^e$ is transported to derived cotensor
over $C^e$.  Therefore
\[
\Tor_\bullet^{A^e}(A,M)\simeq \Cotor_\bullet^{C^e}(C_\Delta,\KK_A^{\mathrm{bi}}(M)),
\]
which is the stated formula.  Taking $M=A$ yields the second assertion.
\end{proof}

\begin{theorem}[Associative chiral Hochschild/coHochschild cohomology; \ClaimStatusProvedHere]
\label{thm:HH-coHH-cohomology}
Under the same hypotheses, for every complete $A$-bimodule $M$ there is a canonical isomorphism
\[
\HH_{\mathrm{ch}}^\bullet(A,M)\simeq
\coHoch_{\mathrm{ch}}^\bullet(C,\KK_A^{\mathrm{bi}}(M)).
\]
In particular,
\[
\HH_{\mathrm{ch}}^\bullet(A)\simeq
\coHoch_{\mathrm{ch}}^\bullet(C)
=
\Coext_{C^e}^\bullet(C_\Delta,C_\Delta).
\]
\end{theorem}

\begin{proof}
Derived Hom is preserved by the equivalence of Theorem~\ref{thm:bimod-bicomod}.  Therefore
\[
\RHom_{A^e}(A,M)\simeq \RHom_{C^e}(C_\Delta,\KK_A^{\mathrm{bi}}(M)),
\]
which is the desired coHochschild cohomology complex.  The specialization to $M=A$ follows from the
diagonal correspondence.
\end{proof}

\begin{corollary}[The annulus as self-cotrace; \ClaimStatusProvedHere]
\label{cor:annulus}
The annular closed object attached to the universal boundary condition $C_\Delta$ is the coHochschild
object
\[
\mathcal H_C^{\mathrm{cl}}
:=
C_\Delta \square_{C^e}^{\mathbf R} C_\Delta
\simeq
\coHoch^{\mathrm{ch}}_\bullet(C)
\simeq
\HH^{\mathrm{ch}}_\bullet(A).
\]
\end{corollary}

\begin{proof}
This is the special case $M=A$ of Theorem~\ref{thm:HH-coHH-homology}.
\end{proof}

\begin{corollary}[Cap action; \ClaimStatusProvedHere]
\label{cor:cap}
There is a canonical action
\[
\coHoch^{\mathrm{ch}}_\bullet(C)\square_C^{\mathbf R} C_\Delta \longrightarrow C_\Delta
\]
obtained by transporting the Hochschild cap action
\[
\HH^{\mathrm{ch}}_\bullet(A)\otimes A\longrightarrow A
\]
across the equivalence of Theorem~\ref{thm:bimod-bicomod}.
\end{corollary}

\begin{proof}
Apply Theorem~\ref{thm:bimod-bicomod} to the classical cap action viewed as a map of bimodules.
\end{proof}

\begin{remark}
The cohomology theorem is the center theorem.  The homology theorem is the annulus theorem.  Together
they say that the algebraic center and algebraic trace of $A$ can be read directly from the
coalgebraic diagonal $C_\Delta$.
\end{remark}

\section{The diagonal bicomodule as universal boundary condition}
\label{sec:open-trace}

There are two operations one wants on the boundary side.

First, interval composition, already realized by derived cotensor and controlled by
Corollary~\ref{cor:tensor-cotensor}.  Second, the ordered pair-of-pants operation, which is not
cotensor but rather the transport of the external tensor product of bimodules.  This second
operation is inherently ordered; in general we do \emph{not} claim it is symmetric without extra
structure.

\begin{definition}[Ordered fusion product]
Let $X,Y\in \Dco_{\mathrm{bi}}(C)$.  Define
\[
X\star Y
:=
\KK_A^{\mathrm{bi}}\Bigl(
\Omega_C^{\mathrm{bi}}(X)\widehat\otimes \Omega_C^{\mathrm{bi}}(Y)
\Bigr),
\]
where $\widehat\otimes$ is the completed tensor product of bimodules with outer $A$-actions:
the left $A$-action is on the first factor and the right $A$-action is on the second factor.
We call $\star$ the \emph{ordered fusion product}.
\end{definition}

\begin{remark}
The adjective ``ordered'' is essential.  Unless further symmetry data are supplied, $\star$ is only
a monoidal product reflecting the ordering of incoming open boundaries.  It is the algebraic shadow
of an ordered pair-of-pants, not yet of the fully symmetric open cobordism category.
\end{remark}

\begin{theorem}[Ordered pair-of-pants algebra; \ClaimStatusProvedHere]
\label{thm:pair-of-pants}
The diagonal bicomodule $C_\Delta$ carries a canonical associative algebra structure with respect to
the ordered fusion product:
\[
\mu_P:C_\Delta\star C_\Delta\longrightarrow C_\Delta,
\qquad
\eta_P:\mathbf 1_\star\longrightarrow C_\Delta,
\]
where $\mathbf 1_\star:=\KK_A^{\mathrm{bi}}(k)$ is the image of the trivial bimodule.
\end{theorem}

\begin{proof}
By Theorem~\ref{thm:diagonal}, $C_\Delta\simeq \KK_A^{\mathrm{bi}}(A)$.  The multiplication and unit
of the algebra $A$,
\[
\mu_A:A\widehat\otimes A\to A,\qquad \eta_A:k\to A,
\]
transport across $\KK_A^{\mathrm{bi}}$ to maps
\[
\mu_P:C_\Delta\star C_\Delta\to C_\Delta,\qquad \eta_P:\mathbf 1_\star\to C_\Delta.
\]
Associativity and unitality follow immediately from associativity and unitality of $A$.
\end{proof}

\begin{definition}[Ordered open trace formalism]
An \emph{ordered genus-zero open trace formalism} consists of:
\begin{enumerate}[label=(\roman*)]
\item a composition product $\square_C^{\mathbf R}$ with unit object $C_\Delta$,
\item an ordered fusion product $\star$,
\item an associative algebra structure $(C_\Delta,\mu_P,\eta_P)$ for $\star$,
\item a closed object $\mathcal H_C^{\mathrm{cl}}$ obtained by annular closure,
\item an annulus action of $\mathcal H_C^{\mathrm{cl}}$ on $C_\Delta$,
\end{enumerate}
subject to the sewing relations induced from the algebra side.
\end{definition}

\begin{theorem}[Ordered genus-zero open trace formalism; \ClaimStatusProvedHere]
\label{thm:ordered-open}
The data
\[
\bigl(\Dco_{\mathrm{bi}}(C),\square_C^{\mathbf R},C_\Delta,\star,\mu_P,\eta_P,\mathcal H_C^{\mathrm{cl}}\bigr)
\]
constitute an ordered genus-zero open trace formalism.  More precisely:
\begin{enumerate}[label=(\alph*)]
\item $C_\Delta$ is the unit for interval composition $\square_C^{\mathbf R}$.
\item $(C_\Delta,\mu_P,\eta_P)$ is an associative algebra object for ordered fusion.
\item Annular closure is
\[
\mathcal H_C^{\mathrm{cl}}\simeq \coHoch^{\mathrm{ch}}_\bullet(C).
\]
\item The annulus acts on the boundary object by the action of
Corollary~\ref{cor:cap}.
\item All ordered genus-zero sewing identities are transported from the corresponding identities for
$A$-bimodules.
\end{enumerate}
\end{theorem}

\begin{proof}
Statement (a) is Corollary~\ref{cor:unit}.  Statement (b) is
Theorem~\ref{thm:pair-of-pants}.  Statement (c) is
Corollary~\ref{cor:annulus}.  Statement (d) is Corollary~\ref{cor:cap}.
Finally, every ordered genus-zero sewing relation on the coalgebra side is the image under the
equivalence $\KK_A^{\mathrm{bi}}$ of the corresponding relation on the algebra side.
\end{proof}

\begin{remark}[What is proved and what is not]
The theorem gives the algebraic ordered open package.  To upgrade it to a fully symmetric
one-brane open TCFT one needs additional structure that symmetrizes ordered fusion---for example a
boundary symmetry datum or a direct geometric realization by bordered configuration spaces.
This is formulated below as part of the conjectural frontier.
\end{remark}

\subsection{Calabi--Yau transport}

The ordered open trace formalism acquires a Frobenius structure under a Calabi--Yau hypothesis on
$A$.

\begin{definition}[Complete chiral Calabi--Yau structure]
A strongly admissible associative chiral algebra $A$ is called \emph{$d$-Calabi--Yau} if:
\begin{enumerate}[label=(\alph*)]
\item $A$ is perfect as an $A^e$-module, and
\item there is an isomorphism in $\Dpbw_{\mathrm{bi}}(A)$
\[
\varphi_A:A\xrightarrow{\sim}\RHom_{A^e}(A,A^e)[d].
\]
\end{enumerate}
Equivalently, $A$ carries a nondegenerate $d$-shifted derived trace pairing.
\end{definition}

\begin{theorem}[Shifted ordered Frobenius structure; \ClaimStatusProvedHere]
\label{thm:CY}
Assume $A$ is $d$-Calabi--Yau.  Then the underlying complex of $C_\Delta$ carries a canonical
$d$-shifted ordered Frobenius structure.  More precisely:
\begin{enumerate}[label=(\roman*)]
\item the ordered pair-of-pants multiplication $\mu_P$ and unit $\eta_P$ are as above;
\item there is a trace
\[
\eps_P:C_\Delta\to k[-d]
\]
transported from the Calabi--Yau trace on $A$;
\item there is a coproduct
\[
\Delta_P:C_\Delta\to C_\Delta\star C_\Delta[d]
\]
adjoint to $\mu_P$ under $\eps_P$;
\item the ordered Frobenius identities
\[
(\mu_P\star \id)\circ(\id\star \Delta_P)
=
\Delta_P\circ \mu_P
=
(\id\star\mu_P)\circ(\Delta_P\star \id)
\]
hold.
\end{enumerate}
\end{theorem}

\begin{proof}
On the algebra side, the Calabi--Yau structure identifies $A$ with its bimodule dual shifted by
$d$.  Therefore multiplication
\[
\mu_A:A\widehat\otimes A\to A
\]
admits a unique adjoint coproduct
\[
\Delta_A:A\to A\widehat\otimes A[d]
\]
with respect to the trace pairing $\eps_A:A\to k[-d]$; explicitly one may define $\Delta_A$ as the
adjoint of $\mu_A$ under $\varphi_A$.  The ordered Frobenius identities are formal consequences of
associativity of $\mu_A$ and the nondegeneracy of the pairing.

Transport the entire package $(\mu_A,\eta_A,\eps_A,\Delta_A)$ across the equivalence
$\KK_A^{\mathrm{bi}}$ and use Theorem~\ref{thm:diagonal}.  The resulting structure maps are
$\mu_P,\eta_P,\eps_P,\Delta_P$, and all identities are preserved because the equivalence preserves
composition.
\end{proof}

\begin{corollary}[Cardy operator on the coalgebra side; \ClaimStatusProvedHere]
\label{cor:cardy}
Under the hypotheses of Theorem~\ref{thm:CY}, the ordered annulus endomorphism of the universal
boundary object is
\[
\mathsf{Ann}_{C_\Delta}:=\mu_P\circ \Delta_P\in \operatorname{End}(C_\Delta)[d].
\]
Its annular closure factors through the closed object
\[
\mathcal H_C^{\mathrm{cl}}\simeq \coHoch^{\mathrm{ch}}_\bullet(C).
\]
\end{corollary}

\begin{proof}
Transport the usual annulus endomorphism $\mu_A\circ \Delta_A$ of the Calabi--Yau algebra $A$ across
$\KK_A^{\mathrm{bi}}$ and use Corollary~\ref{cor:annulus}.
\end{proof}

\section{A master theorem and its conceptual form}
\label{sec:master}

For repository integration it is useful to package the preceding results into a single statement.

\begin{theorem}[Master theorem; \ClaimStatusProvedHere]
\label{thm:master}
Let $A$ be a strongly admissible augmented associative chiral algebra on $X$, assume
$(A,A^{\op})$ is EZ-admissible, and set
\[
C:=\Barch(A).
\]
Then:
\begin{enumerate}[label=(\arabic*)]
\item \textbf{Ordered shuffle.}
There is a quasi-isomorphism of conilpotent dg coalgebras
\[
\Barch(A\chotimes B)\simeq \Barch(A)\widehat\otimes \Barch(B)
\]
for every EZ-admissible partner $B$.

\item \textbf{Opposite-duality.}
There is a canonical isomorphism
\[
\Barch(A^{\op})\simeq \Barch(A)^{\cop}.
\]

\item \textbf{Enveloping duality.}
There is a canonical quasi-isomorphism
\[
\Barch(A^e)\simeq C^e.
\]

\item \textbf{Two-sided bimodule/bicomodule duality.}
There is an exact equivalence
\[
\Dpbw_{\mathrm{bi}}(A)\simeq \Dco_{\mathrm{bi}}(C).
\]

\item \textbf{Diagonal correspondence.}
Under this equivalence one has
\[
{}_A A_A \longleftrightarrow C_\Delta.
\]

\item \textbf{Composition.}
The diagonal bicomodule is the unit for derived cotensor:
\[
C_\Delta\square_C^{\mathbf R} X\simeq X\simeq X\square_C^{\mathbf R}C_\Delta.
\]

\item \textbf{Hochschild/coHochschild homology.}
For every complete $A$-bimodule $M$,
\[
\HH^{\mathrm{ch}}_\bullet(A,M)
\simeq
\coHoch^{\mathrm{ch}}_\bullet(C,\KK_A^{\mathrm{bi}}(M)).
\]

\item \textbf{Hochschild/coHochschild cohomology.}
For every complete $A$-bimodule $M$,
\[
\HH_{\mathrm{ch}}^\bullet(A,M)
\simeq
\coHoch_{\mathrm{ch}}^\bullet(C,\KK_A^{\mathrm{bi}}(M)).
\]

\item \textbf{Tangent theory.}
For a square-zero extension $A_\eps=A\oplus \eps M$,
\[
\Barch(A_\eps)\cong C\oplus \eps\,\KK_A^{\mathrm{bi}}(M)\qquad (\mathrm{mod}\ \eps^2).
\]

\item \textbf{Ordered open trace formalism.}
The coalgebraic boundary object $C_\Delta$ carries ordered pair-of-pants, annulus, and composition
operations transported from the algebra side, and if $A$ is $d$-Calabi--Yau then the underlying
complex of $C_\Delta$ becomes a $d$-shifted ordered Frobenius object.
\end{enumerate}
\end{theorem}

\begin{proof}
This is a summary of Theorems~\ref{thm:shuffle}, \ref{thm:opposite},
\ref{thm:bimod-bicomod}, \ref{thm:diagonal}, \ref{thm:HH-coHH-homology},
\ref{thm:HH-coHH-cohomology}, \ref{thm:tangent=K}, and \ref{thm:ordered-open}, together with
Corollaries~\ref{cor:enveloping}, \ref{cor:unit}, and \ref{cor:infdual}, and
Theorem~\ref{thm:CY}.
\end{proof}

\section{Conjectural frontier}
\label{sec:conjectures}

The preceding sections supply a clean algebraic core.  What remains is to geometricize it and to
relate it to neighboring duality theories.  We now isolate the next statements as explicit
conjectures.

\subsection{Francis--Gaitsgory as the commutator shadow}

The associative story developed here is not the Francis--Gaitsgory Lie/cocommutative story.
Rather, the latter should appear as the commutative shadow of the former.

\begin{conjecture}[Comm\-utator-shadow theorem; \ClaimStatusConjectured]
\label{conj:FG-shadow}
Let $A$ be a filtered strongly admissible associative chiral algebra with exhaustive separated
commutator filtration $F_{\mathrm{com}}^\bullet A$.  Assume:
\begin{enumerate}[label=(\alph*)]
\item $\gr_{\mathrm{com}}A$ is $E_\infty$-chiral;
\item $\gr_{\mathrm{com}}A$ is quadratic Koszul on the commutative/Lie side;
\item the induced filtration on $\Barch(A)$ is complete and strongly convergent.
\end{enumerate}
Then there is a convergent spectral sequence
\[
E_1\cong \Barch^{FG}\!\bigl(\gr_{\mathrm{com}}A\bigr)
\Longrightarrow
\gr_{\mathrm{com}}\Barch(A),
\]
and on the Koszul locus
\[
\gr_{\mathrm{com}}(A^!)\simeq \bigl(\gr_{\mathrm{com}}A\bigr)^!_{FG}.
\]
\end{conjecture}

\begin{remark}[Proof obligations]
A proof should construct the commutator filtration directly on the ordered bar complex, show that the
associated graded kills the non-$\Sigma$ defect and descends to the commutative bar complex, and
identify the induced differential with the Francis--Gaitsgory differential.
\end{remark}

\subsection{Bordered configuration spaces and the geometric open theory}

The algebraic ordered open trace formalism should be realized geometrically by ordered bordered
configuration spaces or, equivalently, by a ribbon/Swiss-cheese type compactification adapted to the
chiral setting.

\begin{conjecture}[Geometric realization of the diagonal bicomodule; \ClaimStatusConjectured]
\label{conj:bordered}
There exists a bordered ordered configuration-space model
\[
\mathfrak C^{\mathrm{bord}}_{g,r,s}(X)
\]
and a functorial system of chain-level operations on the diagonal bicomodule $C_\Delta$ such that:
\begin{enumerate}[label=(\arabic*)]
\item the interval component reproduces the bicomodule $C_\Delta$;
\item gluing of boundary intervals is computed by derived cotensor over $C$;
\item the ordered pair-of-pants operation coincides with $\mu_P$;
\item annular closure computes $\coHoch^{\mathrm{ch}}_\bullet(C)$;
\item the resulting geometric operations recover the ordered open trace formalism of
Theorem~\ref{thm:ordered-open}.
\end{enumerate}
\end{conjecture}

\begin{remark}[Proof obligations]
One needs a bordered Fulton--MacPherson type compactification for ordered points on intervals and
annuli, together with residue or propagator maps realizing the composition and pair-of-pants
operations.  At present the algebraic theory is ahead of the geometry.
\end{remark}

\subsection{Associative modular Maurer--Cartan theory}

The higher-genus modular story should admit an associative analogue built from ordered cyclic
compactifications and the universal boundary object.

\begin{conjecture}[Associative modular Maurer--Cartan class; \ClaimStatusConjectured]
\label{conj:modular}
There exists an associative cyclic deformation complex
\[
\mathrm{Def}^{\mathrm{cyc}}_{\Assch}(A)
\]
and a canonical Maurer--Cartan element
\[
\Theta_A^{\Assch}
\in
\MC\!\Bigl(
\mathrm{Def}^{\mathrm{cyc}}_{\Assch}(A)\widehat\otimes
R\Gamma(\overline{M}_{g,\bullet},k)
\Bigr)
\]
whose linear term is the genus anomaly, whose quadratic term encodes the first pair-of-pants
boundary correction, and whose full Maurer--Cartan equation expresses compatibility of all ordered
higher-genus boundary gluings.
\end{conjecture}

\begin{remark}[Proof obligations]
The missing input is an ordered cyclic compactification whose codimension-one strata reproduce the
combinatorics of boundary splitting and handle attachment.  The relevant signs should be governed by
an ordered analogue of the cyclic bar complex rather than by the symmetric modular operad formalism.
\end{remark}

\subsection{BRST and Drinfeld--Sokolov compatibility}

The last frontier is compatibility between associative self-duality and reduction.

\begin{conjecture}[Reduction commutes with associative chiral duality; \ClaimStatusConjectured]
\label{conj:DS}
Let $Q_{DS}$ denote a BRST/Drinfeld--Sokolov reduction functor defined on a filtered pronilpotent
subcategory of strongly admissible associative chiral algebras.  Assume:
\begin{enumerate}[label=(\alph*)]
\item $Q_{DS}$ is exact on the relevant filtered complexes;
\item $Q_{DS}$ preserves PBW filtrations;
\item the BRST differential lifts to the ordered bar coalgebra.
\end{enumerate}
Then
\[
Q_{DS}(A^!)\simeq Q_{DS}(A)^!.
\]
Moreover, for every $A$-bimodule $M$ one should have
\[
Q_{DS}\bigl(\KK_A^{\mathrm{bi}}(M)\bigr)\simeq
\KK_{Q_{DS}(A)}^{\mathrm{bi}}\bigl(Q_{DS}(M)\bigr).
\]
\end{conjecture}

\begin{remark}
If true, this would show that associative chiral self-duality commutes not only with passage from
algebras to coalgebras but also with reduction to the corresponding $W$-type objects.
\end{remark}

\section{Repository integration notes}
\label{sec:integration}

For direct integration into a research repository, the following replacements are recommended.

\begin{enumerate}[label=(\arabic*)]
\item Replace the abstract hypotheses in \S\ref{sec:setup} by the precise labels of the repository's
existing bar--cobar and module/comodule theorems.

\item Replace the shorthand notation $\widehat\otimes$ and $\chotimes$ by the repository's preferred
completed tensor notation and chiral monoidal symbol.

\item If the repository already has a sign appendix for the ordered bar complex, point the proofs of
Theorems~\ref{thm:opposite} and \ref{thm:tangent=K} to that appendix.

\item If the repository distinguishes coderived and contraderived categories more sharply, split
Theorem~\ref{thm:bimod-bicomod} accordingly.

\item If a bordered configuration-space chapter is later added, Conjecture~\ref{conj:bordered}
should be moved forward and rephrased as a theorem package with explicit chain models.
\end{enumerate}

\appendix

\section{Appendix: signs and conventions}

All formulas above are invariant under the usual choice of suspension conventions, provided one keeps
the same convention simultaneously in the ordered bar complex, the shuffle map, and the twisting
differential.  For clarity we summarize the minimal sign content.

\begin{enumerate}[label=(\arabic*)]
\item Reversal in Theorem~\ref{thm:opposite} uses the standard suspended Koszul sign
\[
(-1)^{\sum_{i<j}(|a_i|+1)(|a_j|+1)}.
\]

\item The left twisting term in Definition~\ref{def:Kbi} carries the sign needed to move the
coderivation on $C$ past the first coalgebra leg before applying the left $A$-action.

\item The right twisting term carries the additional sign needed to move the twisting cochain across
the middle bimodule factor.

\item In the strongly admissible regime the filtration is complete and bounded below, so all spectral
sequence and homotopy-transfer arguments used in the proofs converge.
\end{enumerate}

\section{Appendix: new results proved in this note}

For quick reference, here is the list of new results proved here rather than taken as repository
inputs.

\begin{enumerate}[label=(\arabic*)]
\item Ordered chiral shuffle theorem (Theorem~\ref{thm:shuffle}).
\item Opposite-duality and anti-involution transport (Theorem~\ref{thm:opposite},
Corollary~\ref{cor:anti}).
\item Enveloping duality (Corollary~\ref{cor:enveloping}).
\item Two-sided twisted bicomodule and its dg structure (Definition~\ref{def:Kbi},
Lemma~\ref{lem:Kbi-dg}).
\item Identification of the tangent ordered bar complex with the twisted bicomodule
(Theorem~\ref{thm:tangent=K}).
\item PBW-complete bimodule/bicomodule equivalence (Theorem~\ref{thm:bimod-bicomod}).
\item Diagonal correspondence (Theorem~\ref{thm:diagonal}).
\item Composition unit theorem (Corollary~\ref{cor:unit}).
\item Tensor--cotensor gluing (Corollary~\ref{cor:tensor-cotensor}).
\item Hochschild/coHochschild homology and cohomology equivalence
(Theorems~\ref{thm:HH-coHH-homology} and \ref{thm:HH-coHH-cohomology}).
\item Infinitesimal annular variation (Proposition~\ref{prop:infann}).
\item Ordered genus-zero open trace formalism (Theorem~\ref{thm:ordered-open}).
\item Calabi--Yau transport to an ordered Frobenius structure (Theorem~\ref{thm:CY}).
\end{enumerate}

\bigskip

\noindent\textbf{Final principle.}
The diagonal is the boundary.  The bar coalgebra is the ambient ordered phase space.  The annulus is
the self-cotensor of the diagonal.  Once this is recognized, the associative chiral duality ceases
to be merely an equivalence between algebra and coalgebra and becomes a machine for open traces,
centers, and eventually bordered genus.

